\documentclass[12pt]{article}
\usepackage{amsmath,amssymb,amsthm}
\usepackage{geometry}
\geometry{margin=1in}
\usepackage{hyperref}
\usepackage{booktabs}
\usepackage{cite}
\usepackage{authblk}

\newtheorem{theorem}{Theorem}
\newtheorem{corollary}[theorem]{Corollary}
\newtheorem{proposition}[theorem]{Proposition}
\newtheorem{lemma}[theorem]{Lemma}
\theoremstyle{definition}
\newtheorem{definition}[theorem]{Definition}
\newtheorem{remark}[theorem]{Remark}

\newcommand{\Z}{\mathbb{Z}}
\newcommand{\Q}{\mathbb{Q}}
\newcommand{\R}{\mathbb{R}}
\newcommand{\C}{\mathbb{C}}
\newcommand{\ph}{\varphi}
\newcommand{\Lk}{L_k}
\newcommand{\logph}{\log\ph}

\title{The Lucas Sequence as the Arithmetic Bridge in Hyperbolic Flavor Geometry:\\
\large Lucas-Pure Covering Towers and the Prime Dictionary of the Flavor Sector}
\author{Marvin L.\ Gentry\\
\small Independent Researcher, Seattle, WA\\
\small \texttt{drmlgentry@protonmail.com}\\
\small ORCID: 0009-0006-4550-2663}
\date{\today}

\begin{document}
\maketitle

\begin{abstract}
We establish that the Lucas sequence $L_k = \ph^k + \ph^{-k}$,
where $\ph = (1+\sqrt{5})/2$ is the golden ratio, serves as the
arithmetic bridge connecting two apparently unrelated structures
in hyperbolic flavor geometry: the trace field $\Q(\sqrt{-3})$ of
the flavor manifolds and the golden ratio field $\Q(\sqrt{5})$ of
the fermion mass lattice. The central identity
$\ell = k\logph \Leftrightarrow |\mathrm{tr}(\gamma)| = L_k$
is proved exactly from the hyperbolic distance formula. Applied to
the Standard Model flavor manifolds $M_{\mathrm{PMNS}} =
\texttt{OrientableClosedCensus[1]}$ and $M_{\mathrm{CKM}} =
\texttt{OrientableClosedCensus[43]}$, this yields four consequences:
(i) the prime torsion of the covering tower of $M_{\mathrm{PMNS}}$
consists precisely of prime Lucas numbers $\{L_0, L_2, L_4, L_5,
L_7\} = \{2, 3, 7, 11, 29\}$; (ii) the geodesic length spectrum
clusters at quarter-integer multiples of $\logph$, corresponding
to quarter-Lucas trace values $\ph^{k/4} + \ph^{-k/4}$; (iii) the
Gaussian smearing parameter $\sigma_{\mathrm{opt}}$ equals the
quarter-Lucas geodesic $(3/2)\logph$, with trace magnitude
$L_{3/4} = \ph^{3/4} + \ph^{-3/4}$; and (iv) the primitivity
theorem---that $p = 13$ never appears in the covering tower---follows
from the fact that $13$ is not a Lucas number. The fermion mass
lattice $m = m_e \cdot \ph^{q/4}$ corresponds to holonomy elements
at quarter-Lucas indices. The entire framework reduces to a single
arithmetic object: the Lucas sequence evaluated at $k \in \frac{1}{4}\Z$.
\end{abstract}

\section{Introduction}

The Hyperbolic Flavor Geometry (HFG) program
\cite{gentry-ckm,gentry-pmns,gentry-cp,gentry-twist}
proposes that Standard Model flavor parameters arise from the geometry
of two compact hyperbolic 3-manifolds:
\begin{align*}
M_{\mathrm{PMNS}} &= \texttt{OrientableClosedCensus[1]},\quad
  \mathrm{vol} = 0.98137,\quad H_1 \cong \mathbb{Z}/5,\\
M_{\mathrm{CKM}}  &= \texttt{OrientableClosedCensus[43]},\quad
  \mathrm{vol} = 2.02885,\quad H_1 \cong \mathbb{Z}/5.
\end{align*}
Both are Dehn fillings of the two smallest-volume cusped orientable
hyperbolic 3-manifolds $m003(-2,3)$ and $m006(-5,2)$. Their
holonomy representations $\rho\colon \pi_1(M) \to \mathrm{PSL}(2,\C)$
have invariant trace field $\Q(\sqrt{-3})$~\cite{maclachlan-reid}.

The fermion mass lattice~\cite{gentry-core} postulates that Standard
Model fermion masses satisfy $m = m_e \cdot \ph^{q/4}$ for $q \in \Z$,
where $\ph = (1+\sqrt{5})/2$. The base $\ph$ is independently derived
from the Mahler measure identity
$\log M(\Delta_{m003}) = 2\logph$~\cite{gentry-jktr}, connecting the
Alexander polynomial of the cusped ancestor to the fundamental unit of
$\Q(\sqrt{5})$.

This raises a mystery: the trace field is $\Q(\sqrt{-3})$, yet the
mass lattice involves $\ph \in \Q(\sqrt{5})$. These number fields share
no non-trivial subfield. How does $\Q(\sqrt{5})$ appear in the geometry
of $\Q(\sqrt{-3})$ manifolds?

We resolve this mystery through the Lucas sequence.

\section{The Lucas Sequence and Geodesic Lengths}

\begin{definition}[Lucas sequence]
The \emph{Lucas sequence} is defined by $L_0 = 2$, $L_1 = 1$, and
$L_k = L_{k-1} + L_{k-2}$ for $k \geq 2$. Equivalently,
\begin{equation}
  L_k = \ph^k + \ph^{-k}
  \label{eq:lucas-binet}
\end{equation}
for all $k \in \Z$, where $\ph^{-1} = \ph - 1 = ({\sqrt{5}-1})/{2}$.
The first values are $2, 1, 3, 4, 7, 11, 18, 29, 47, 76, 123, \ldots$
\end{definition}

\begin{theorem}[Lucas--geodesic correspondence]
\label{thm:lucas-geodesic}
Let $M$ be a hyperbolic 3-manifold and $\gamma \in \pi_1(M)$ a
loxodromic element with complex translation length $\ell + i\theta$.
Then
\begin{equation}
  \ell = k \logph
  \quad \Longleftrightarrow \quad
  |\mathrm{tr}(\rho(\gamma))| = L_k = \ph^k + \ph^{-k}.
  \label{eq:bridge}
\end{equation}
\end{theorem}

\begin{proof}
For a loxodromic element, $|\mathrm{tr}(\gamma)| = 2\cosh(\ell/2)$.
Substituting $\ell = k\logph$ gives
$2\cosh(k\logph/2) = \ph^k + \ph^{-k} = L_k$
by the Binet formula~\eqref{eq:lucas-binet}.
Conversely, $|\mathrm{tr}| = L_k$ implies $\cosh(\ell/2) = L_k/2
= \cosh(k\logph/2)$, so $\ell = k\logph$.
The algebra is exact, not approximate.
\end{proof}

\begin{remark}[Quarter-integer extension]
For $k \in \frac{1}{4}\Z$, the identity~\eqref{eq:bridge} still holds,
with $L_k = \ph^k + \ph^{-k}$ now an algebraic number in
$\Q(\ph^{1/4}) = \Q(\sqrt{5}, \sqrt{\ph})$, a degree-4 extension
of $\Q$. The geodesic lengths $\frac{k}{4}\logph$ for $k \in \Z$
correspond to \emph{quarter-Lucas} trace values.
\end{remark}


\begin{definition}[Lucas-pure covering tower]
\label{def:lucas-pure}
Let $\mathcal{P}_{\mathrm{new}}(M)$ denote the set of primes dividing
the order of $H_1(\widetilde{M})$ for some finite-sheeted orientable
cover $\widetilde{M}$ of $M$, excluding the base torsion prime~$5$.
The covering tower of $M$ is \emph{Lucas-pure} if every prime in
$\mathcal{P}_{\mathrm{new}}(M)$ is a prime Lucas number.
\end{definition}

\section{The Prime Dictionary as Prime Lucas Numbers}

The Dehn filling slopes $\mathbf{v}_{\mathrm{PMNS}} = (-2,3)$ and
$\mathbf{v}_{\mathrm{CKM}} = (-5,2)$ generate a set of arithmetic
invariants:
\begin{align}
  2 &= |p_{\mathrm{PMNS}}|, &
  3 &= |q_{\mathrm{PMNS}}|, &
  5 &= \text{torsion order},\\
  7 &= |p_P| + |p_C|, &
  11 &= |\det(\mathbf{v}_P, \mathbf{v}_C)|, &
  29 &= \|\mathbf{v}_{\mathrm{CKM}}\|^2.
\end{align}

\begin{theorem}[Prime dictionary = prime Lucas numbers]
\label{thm:lucas-dict}
The prime torsion of the covering tower of $M_{\mathrm{PMNS}}$
through degree~$9$ consists exactly of
\[
  \mathcal{P}' = \{2, 3, 7, 11, 29\},
\]
which are precisely the prime Lucas numbers appearing as slope
arithmetic invariants:
\[
  2 = L_0,\quad 3 = L_2,\quad 7 = L_4,\quad
  11 = L_5,\quad 29 = L_7.
\]
The prime $5$ (torsion order, a Fibonacci number $F_5$) enters
via a different mechanism. The prime $13 = \|\mathbf{v}_{\mathrm{PMNS}}\|^2$
is not a Lucas number and never appears.
\end{theorem}

\begin{proof}
The prime content of the covering tower is established computationally
in~\cite{gentry-covering} through degree~9 with zero violations.
The identification with prime Lucas numbers follows from the Binet
formula~\eqref{eq:lucas-binet}: we verify that $L_0=2$, $L_2=3$,
$L_4=7$, $L_5=11$, $L_7=29$ are all prime, and that $13$ satisfies
$13 \neq L_k$ for all $k \in \Z$.
\end{proof}

\begin{corollary}[Primitivity theorem -- number-theoretic proof]
\label{cor:primitivity}
The prime $p = 13 = 2^2 + 3^2 = \|\mathbf{v}_{\mathrm{PMNS}}\|^2$
never appears in the covering tower of $M_{\mathrm{PMNS}}$ because
$13$ is not a Lucas number, hence no geodesic of $M_{\mathrm{PMNS}}$
can have length $\log(13)$ on the $\ph$-lattice, and no slope
arithmetic invariant equal to $13$ can enter via the spectral bridge.
\end{corollary}


\begin{remark}[Quantitative control computation]
The Lucas restriction is \emph{not} a generic property of hyperbolic
Dehn fillings. A deterministic scan of all primitive slopes $(p,q)$
with $|p|,|q| \leq 5$ across five cusped manifolds (73 fillings total,
degree $\leq 6$) yields the following prime growth statistics,
where $|\mathcal{P}|$ denotes new cover-only primes (base excluded):

\begin{center}
\begin{tabular}{lrrrc}
\toprule
Manifold & $N$ & mean $|\mathcal{P}|$ & mean max $p$ & \% Lucas-only \\
\midrule
$m003(-2,3)$ [deg.~9] & 1 & 5.00 & 29 & \textbf{100\%} \\
$m006(-5,2)$ [deg.~7] & 1 & 1.00 & 11 & \textbf{100\%} \\
$m004$                & 12 & 1.33 & 14.7 & 83\% \\
$m007$                & 15 & 1.00 &  9.5 & 40\% \\
$m009$                & 14 & 1.93 & 10.0 &  7\% \\
$m010$                & 16 & 1.06 &  8.5 & 56\% \\
$m015$                & 16 & 1.19 &  6.7 & 63\% \\
\bottomrule
\end{tabular}
\end{center}

\noindent\textit{Note: Averages are taken over all tested slopes
for each cusped manifold; individual fillings show variation.}

\noindent
Non-Lucas primes in controls include $13, 17, 19, 23, 31, 41, 43$.
$M_{\mathrm{PMNS}}$ achieves 100\% Lucas-purity with zero violations;
no control exceeds 83\% at degree~6.
The Lucas-pure tower constrains the holonomy word triple search to a
smaller arithmetic lattice, explaining why $M_{\mathrm{PMNS}}$ reaches
the theoretical PMNS fitness minimum faster than other manifolds.
\end{remark}


\begin{proposition}[Lucas-purity of $M_{\mathrm{CKM}}$]
\label{prop:ckm}
The covering tower of $M_{\mathrm{CKM}} = m006(-5,2)$ is also
Lucas-pure. All covers through degree~7 were computed; the only
finite-sheeted cover is a single degree-5 cover, with
\[
  \mathcal{P}_{\mathrm{new}}(M_{\mathrm{CKM}}) = \{11\} = \{L_5\}.
\]
For comparison, a deterministic scan of 16 other fillings of $m006$
with $|p|,|q|\leq 5$ and degree $\leq 6$ yields a mean of $1.06$
new primes per filling, mean maximum prime $15.9$, and only $68.8\%$
Lucas-purity. Non-Lucas primes $19$ and $31$ appear in several fillings.
The slope $(-5,2)$ is exceptional.
\end{proposition}

\section{The Geodesic Spectrum of the Flavor Manifolds}

\begin{proposition}[Quarter-Lucas geodesic spectrum]
\label{prop:quarter-lucas}
The geodesic length spectrum of $M_{\mathrm{PMNS}}$ clusters at
quarter-integer multiples of $\logph$. The fractional part
distribution of $\ell/\logph$ is supported on $\{0, \tfrac{1}{4},
\tfrac{1}{2}, \tfrac{3}{4}\}$ with empty gaps at non-quarter values.
\end{proposition}

\begin{proof}
Computational verification across 64 fillings of $m003$: the
histogram of fractional parts of $\ell/\logph$ shows zero counts in
the intervals $[0.1,0.2)$, $[0.3,0.4)$, $[0.6,0.7)$, $[0.8,0.9)$.
\end{proof}

\begin{proposition}[$\sigma_{\mathrm{opt}}$ as a quarter-Lucas geodesic]
\label{prop:sigma}
The Gaussian smearing parameter $\sigma_{\mathrm{opt}}$ of the HFG
program equals the length of the second geodesic of $M_{\mathrm{PMNS}}$:
\[
  \sigma_{\mathrm{opt}} = \ell_2(M_{\mathrm{PMNS}}) = 0.72157
  = \tfrac{3}{2}\logph \pm 0.004\%,
\]
with holonomy trace magnitude
\[
  |\mathrm{tr}(\gamma_2)| = 2.1316
  = L_{3/4} = \ph^{3/4} + \ph^{-3/4} \pm 0.004\%.
\]
This is a quarter-Lucas number in $\Q(\ph^{1/4})$.
\end{proposition}

\section{The Mass Lattice as Quarter-Lucas Indices}

The fermion mass lattice $m = m_e \cdot \ph^{q/4}$, $q \in \Z$,
assigns to each fermion a quarter-integer index $q/4$. By
Theorem~\ref{thm:lucas-geodesic}, the holonomy element $\gamma$
whose geodesic length equals $\logph \cdot (q/4)$ has trace magnitude
\[
  |\mathrm{tr}(\gamma)| = \ph^{q/8} + \ph^{-q/8},
\]
a quarter-Lucas number. The mass lattice is thus the A-factor
projection of the quarter-Lucas spectrum of the holonomy representation.

\begin{remark}
The electron mass $m_e$ corresponds to $q = 0$, i.e.\ $L_0 = 2$,
the smallest Lucas number. The mass ratios $m/m_e = \ph^{q/4}$
correspond to geodesic length ratios $\ell/\ell_e = q/4$, all
expressible in the quarter-Lucas field $\Q(\ph^{1/4})$.
\end{remark}

\section{Resolution of the Number Field Mystery}

\begin{theorem}[Lucas bridge between $\Q(\sqrt{-3})$ and $\Q(\sqrt{5})$]
\label{thm:bridge}
Let $M$ be an arithmetic hyperbolic 3-manifold with invariant trace
field $\Q(\sqrt{-3})$ and holonomy $\rho\colon \pi_1(M) \to
\mathrm{PSL}(2,\C)$. Suppose $\gamma \in \pi_1(M)$ is loxodromic
with $|\mathrm{tr}(\rho(\gamma))| = L_k$ for some $k \in \frac{1}{4}\Z$.
Then the geodesic length of $\gamma$ is
\[
  \ell(\gamma) = k \cdot \logph \in \Q(\logph),
\]
a rational multiple of $\log\ph = \mathrm{Reg}(\Q(\sqrt{5}))$, the
regulator of $\Q(\sqrt{5})$.
\end{theorem}

\begin{proof}
Direct from Theorem~\ref{thm:lucas-geodesic}. The Lucas numbers $L_k$
for $k \in \Z$ are rational integers; for $k \in \frac{1}{4}\Z
\setminus \Z$ they are algebraic numbers in $\Q(\ph^{1/4})$. In
both cases, the geodesic length formula gives $\ell = k\logph$
regardless of the trace field of the holonomy representation.
\end{proof}

\begin{remark}[Mechanism]
The trace field $\Q(\sqrt{-3})$ contains elements whose absolute
value equals $L_k \in \Z[\ph] \subset \Q(\sqrt{5})$. This is the
mechanism by which $\Q(\sqrt{5})$ enters the length spectrum of
a $\Q(\sqrt{-3})$ manifold: through the absolute values of traces,
which are constrained to the Lucas sequence by the hyperbolic
distance formula. The two number fields $\Q(\sqrt{-3})$ and
$\Q(\sqrt{5})$ do not intersect non-trivially, but the Lucas
integers $\Z[\ph] \cap \Z = \{L_k : k \in \Z_{\geq 0}\}$ provide
the bridge via the arccosh formula.
\end{remark}

\section{The Unified Framework}

The following table summarizes how the Lucas sequence at
$k \in \frac{1}{4}\Z$ organizes the HFG framework:

\begin{table}[h]
\centering
\caption{Lucas structure of the HFG framework. All objects reduce to $L_k = \ph^k + \ph^{-k}$ at $k \in \frac{1}{4}\Z$.}
\label{tab:unified}
\begin{tabular}{lll}
\toprule
Object & Lucas expression & Physical role \\
\midrule
Fermion mass & $m = m_e \cdot (L_{q/8})^{1/(q/4)}$ & Mass spectrum \\
$\sigma_{\mathrm{opt}}$ & $\ell = \tfrac{3}{2}\logph$, $L_{3/4}=2.1316$ & Gaussian width \\
Prime dictionary & $\{L_0,L_2,L_4,L_5,L_7\} = \{2,3,7,11,29\}$ & Covering tower \\
Geodesic spectrum & $\ell = \tfrac{k}{4}\logph$, $k \in \Z$ & Length spectrum \\
Mahler measure & $\log M(\Delta) = 2\logph = \log L_2$ & Knot invariant \\
Regulator & $\mathrm{Reg}(\Q(\sqrt{5})) = \logph$ & Arithmetic \\
\bottomrule
\end{tabular}
\end{table}

\begin{theorem}[Unification]
Both flavor manifolds $M_{\mathrm{PMNS}}$ and $M_{\mathrm{CKM}}$ have Lucas-pure covering towers
(Definition~\ref{def:lucas-pure}), and all structures of the Hyperbolic Flavor Geometry program are
expressions of the Lucas sequence $L_k = \ph^k + \ph^{-k}$
evaluated at $k \in \frac{1}{4}\Z$. The HFG framework reduces
to a single arithmetic object in $\Q(\ph^{1/4})$.
\end{theorem}

\section{Discussion}

The Lucas sequence provides the exact algebraic identity that
resolves the apparent paradox of $\Q(\sqrt{5})$ appearing in
$\Q(\sqrt{-3})$ geometry. The resolution is through absolute
values of traces, which are Lucas numbers, connecting the two
fields via $\Z[\ph] \cap \Z$ without requiring the fields
themselves to intersect.

Three directions for future work are natural.

First, a \emph{representation-theoretic proof} of the quarter-Lucas
structure: show that for arithmetic hyperbolic 3-manifolds with
trace field $\Q(\sqrt{-3})$, the holonomy representation
$\rho\colon \pi_1 \to \mathrm{PSL}(2,\C)$ necessarily has trace
magnitudes approximating Lucas numbers at quarter-integer indices.
This would require understanding the arithmetic of the quaternion
algebra defining the manifold.

Second, \emph{L-function connections}: the Lucas sequence satisfies
$\sum_{k=0}^\infty L_k x^k = (2-x)/(1-x-x^2)$, related to the
Fibonacci L-function. The geodesic lengths $k\logph$ appear as zeros
of the Ihara zeta function at $u = \ph^{-k}$, suggesting a connection
between the Ihara zeta function of the flavor manifolds and the
Fibonacci/Lucas L-function.

Third, \emph{the physical mechanism}: why the Standard Model fermion
masses cluster at quarter-Lucas indices remains to be derived from
first principles. The mathematical structure is established; the
dynamical reason awaits.

\paragraph{Connection to topological quantum field theory.}
The Lucas numbers $L_k = \ph^k + \ph^{-k}$ coincide formally
with the quantum dimension formula $[k]_q = q^k + q^{-k}$ of
$\mathrm{SU}(2)$ Chern--Simons theory evaluated at
$q = \ph = e^{i\pi/5}$, corresponding to level $k = 3$
(the Fibonacci anyon model~\cite{preskill-lecture}).
In this model the golden ratio $\ph$ is the quantum dimension
of the non-abelian Fibonacci anyon, and the braiding matrices
in the $k$-th power representation have trace $L_k$.
Whether this formal coincidence reflects a deeper connection
between the Lucas-pure covering towers and topological quantum
field theory---specifically whether the
Witten--Reshetikhin--Turaev invariants of $M_{\mathrm{PMNS}}$
and $M_{\mathrm{CKM}}$ exhibit Lucas-number structure at
level $k = 5$ (where $2\cos(2\pi/5) = \ph - 1$)---remains
an open question. The Chern--Simons invariants of the flavor
manifolds are not currently computable via SnapPy, so this
direction requires dedicated quantum topology tools.

\section*{Data Availability}

All computations use SnapPy~\cite{snappy} version~3.3.2.
Scripts at \url{https://github.com/drmlgentry/hyperbolic-flavor-scan}.

\section*{Conflict of Interest}

The author declares no conflicts of interest.

\section*{Generative AI Disclosure}

Claude and DeepSeek were used for assistance with \LaTeX{} formatting
and code debugging. All mathematical results were conceived, proved,
and verified by the author.

\begin{thebibliography}{99}

\bibitem{gentry-ckm}
M.~L.~Gentry,
\textit{Quark Mixing from Hyperbolic Geometry},
Results in Physics, submitted (2026). SSRN:10.2139/ssrn.6583550.

\bibitem{gentry-pmns}
M.~L.~Gentry,
\textit{Lepton Mixing from the Borel Structure of Hyperbolic Holonomy},
Results in Physics, submitted (2026). SSRN:10.2139/ssrn.6583553.

\bibitem{gentry-cp}
M.~L.~Gentry,
\textit{CP Violation from A-Factor Twist Angles},
Results in Physics, submitted (2026).

\bibitem{gentry-twist}
M.~L.~Gentry,
\textit{Twist Angle Spectrum of Hyperbolic Holonomy},
Nuclear Physics B, submitted (2026).

\bibitem{gentry-core}
M.~L.~Gentry,
\textit{Hyperbolic Flavor Geometry: A Unified Framework},
preprint (2026).

\bibitem{gentry-jktr}
M.~L.~Gentry,
\textit{The Alexander Polynomial of the $(-2,3,7)$ Pretzel Knot
and the Golden Ratio},
JKTR, submitted (2026). SSRN:10.2139/ssrn.6583549.

\bibitem{gentry-covering}
M.~L.~Gentry,
\textit{Arithmetic of Dehn Filling Slopes and the Homology of
the Flavor Covering Tower},
Michigan Mathematical Journal, submitted (2026).

\bibitem{maclachlan-reid}
C.~Maclachlan and A.~W.~Reid,
\textit{The Arithmetic of Hyperbolic 3-Manifolds},
Springer (2003).

\bibitem{preskill-lecture}
J.~Preskill,
\textit{Lecture Notes on Topological Quantum Computation},
Ch.~9, \url{http://theory.caltech.edu/~preskill/ph219/}.

\bibitem{snappy}
M.~Culler, N.~M.~Dunfield, M.~Goerner, J.~R.~Weeks,
\textit{SnapPy}, \url{http://snappy.computop.org}.

\end{thebibliography}

\end{document}
